\section{Chromatic and equivariant preliminaries}\label{sec:background}

We work \(p\)-locally. The category \(\Sp^G\) is the stable category of genuine \(G\)-spectra on a complete universe; its compact objects are retracts of \(p\)-local finite genuine \(G\)-CW spectra. All limits, tensor products, and fixed-point functors are derived. When a localized category is used, its tensor product is the corresponding localized smash product.

\subsection{Chromatic tests and finite localization}

For \(i>0\), choose a finite type-\(i\) spectrum with a periodic self-map and write \(T(i)\) for its telescope. Put \(T(0)=K(0)=H\Q\). The Bousfield class of \(T(i)\) is independent of the choice. For a finite spectrum \(V\),
\begin{equation}\label{eq:finite-tests}
 T(i)\wedge V=0\quad\Longleftrightarrow\quad K(i)\wedge V=0.
\end{equation}
The type of a nonzero finite spectrum is its least nonvanishing Morava height. Nilpotence and periodicity give these facts and finite spectra of every type \cite{HS}.

Write \(L_m^f\) for localization at \(T(0)\vee\cdots\vee T(m)\). This is a smashing finite localization; its acyclic compact ideal is generated by a finite spectrum of type \(m+1\). Consequently an \(L_m^f\)-local spectrum is zero if and only if its smashes with \(T(0),\ldots,T(m)\) vanish. We use the telescopic fracture square, a homotopy pullback
\begin{equation}\label{eq:fracture}
\begin{tikzcd}[column sep=large,row sep=large]
 L_m^fD \arrow[r]\arrow[d]&L_{T(m)}D\arrow[d]\\
 L_{m-1}^fD\arrow[r]&L_{m-1}^fL_{T(m)}D
\end{tikzcd}
\end{equation}
for \(m>0\) and \(D\) already \(L_m^f\)-local; equivalently take \(D=L_m^fS\) for the sphere square used below. The square follows from the complementary localizations at the top telescope and the lower finite heights: \(L_{T(m)}L_{m-1}^f=0\), and the two localizations jointly detect \(L_m^f\)-local equivalences.

\subsection{Genuine fixed points and cofree spectra}

Inflation of an ordinary spectrum \(D\) will be denoted by \(\operatorname{inf}_G D\). It satisfies
\[
 \Phi^H(\operatorname{inf}_G D)\simeq D
 \qquad(H\le G).
\]
Geometric fixed points are exact and symmetric monoidal, preserve finite objects, and jointly detect equivalences when taken over all subgroups. For a normal subgroup, their genuine quotient versions are transitive.

The functor from genuine \(P\)-spectra to underlying \(P\)-diagrams has a cofree right adjoint. For the constant diagram with value \(D\), put
\[
 \Bor_P(D)=F(EP_+,\operatorname{inf}_P D),
 \qquad U_P(D)=\Phi^P\Bor_P(D).
\]
We also use \(\Bor_P\) for that right adjoint on a general \(P\)-diagram, specifying the diagram when its action is nontrivial. Its categorical fixed points on a constant diagram are \(D^{BP}=F(BP_+,D)\). The functor \(U_P\) is exact and lax symmetric monoidal. In particular it takes commutative algebras and their modules to commutative algebras and their modules.

If \(X\) is dualizable, adjunction gives the cofree projection formula
\begin{equation}\label{eq:cofree-projection}
 X\otimes\Bor_P(M)\simeq
 \Bor_P\bigl(\operatorname{res}X\otimes M\bigr),
\end{equation}
where \(\operatorname{res}X\) retains its underlying \(P\)-action. We use the usual restriction, coinduction, and normal geometric-localization formulas for genuine spectra \cite[Section~2]{BS}.

\begin{proposition}[Cyclic upper bound]\label{prop:cyclic-upper}
For every finite \(p\)-group \(G\), subgroup \(H\), and \(n\ge0\),
\(r_n(G,H)\le\ell(G,H)\).
\end{proposition}
\begin{proof}
The finite abelian chromatic Smith theorem gives loss at most one for every nontrivial cyclic \(p\)-group, including target height zero \cite[Theorem~2.1 and Corollary~3.12]{BHNNNS}; see also \cite[Theorem~5.1]{KL}. Along a length-\(s\) cyclic subnormal chain, apply it to the finite genuine \(G_i/G_{i-1}\)-spectrum obtained by restricting \(X\) to \(G_i\) and taking genuine geometric \(G_{i-1}\)-fixed points. Transitivity identifies the underlying and full fixed points with \(\Phi^{G_{i-1}}X\) and \(\Phi^{G_i}X\). Iterating at heights \(n+s-1,\ldots,n\) proves the assertion.
\end{proof}

\subsection{Ambidexterity and Burnside marks}

The relevant ambidexterity belongs to the development of chromatic
Tate vanishing and higher semiadditivity. Kuhn proved telescopic
finite-group Tate vanishing \cite[Theorem~1.5]{KuhnTate}; Hopkins
and Lurie established higher ambidexterity in the Morava-local
category \cite[Theorem~0.0.2]{HL}. The telescopic form below is due
to Carmeli, Schlank, and Yanovski. None of these uses identifies
telescopic localization with Morava localization on arbitrary spectra.

We use the published ambidexterity theorem for \(\Sp_{T(l)}\), \(\Sp_{K(l)}\), and \(K(l)\)-local \(E_l\)-modules \cite[Theorem~5.3.1 and Corollary~5.3.2]{CSY}. In these categories every finite groupoid is ambidextrous. Integration along maps of finite groupoids satisfies composition, homotopy base change, and the projection formula; for finite coverings it is the ordinary sum transfer \cite[Section~3.1]{CSY}. Cochains on a finite groupoid with values in the unit are dualizable \cite[Proposition~3.3.7]{CSY}. Tensoring preserves finite-groupoid limits, since it preserves their ambidextrous colimits \cite[Corollary~3.2.4]{CSY}. Thus the canonical exterior product of unit-cochains is an algebra equivalence
\[
 \one^{BY}\otimes\one^{BZ}\simeq\one^{B(Y\times Z)}
\]
for finite groups \(Y,Z\), and similarly for finite groupoids.

For a finite group \(F\), its rational Burnside ring is the product of copies of \(\Q\) indexed by conjugacy classes of subgroups. The \(I\)-mark of a finite \(F\)-set \(X\) is \(|X^I|\). We denote by \(e_I^F\) the idempotent with mark one at the conjugacy class of \(I\) and zero at all others. A finite \(F\)-set \(X\) acts on cochains by multiplication by the transfer of \(1\) along \(X\mathbin{//}F\to BF\); for an orbit, this is restriction followed by transfer. This agrees with the genuine Burnside action. The full mark \(e_F^F\) will be used to isolate full geometric fixed points.

\subsection{Prime points and finite truncation}

For the rest of the paper, primes are indexed by the \emph{tested height}:
\[
 \mathcal P_i^G(H)=
 \{X\in(\Sp^G)^c:K(i)_*(\Phi^H X)=0\},
 \qquad i\ge0.
\]
Thus \(\mathcal P_i^G(H)\) has index \(i+1\) in the convention of \cite{BS}. We suppress \(G\) when it is fixed. The support of a compact object \(X\) consists of the primes not containing it. In particular
\begin{equation}\label{eq:prime-closure}
 q\in\overline{\{y\}}\quad\Longleftrightarrow\quad q\subseteq y.
\end{equation}

\begin{proposition}[Uniform finite localization]\label{prop:finite-local}
Let \(\mathcal C_m^G\) be the category of genuine \(G\)-spectra localized by the inflated sphere \(L_m^fS\). It is rigidly compactly generated. Its compact tensor-triangular spectrum is the finite subspace
\[
 \{(H,i):H\le G\text{ up to conjugacy},\ 0\le i\le m\}
\]
of the equivariant spectrum, with \((H,i)\) representing \(\mathcal P_i^G(H)\). The image of a finite genuine spectrum is compact and dualizable. An object \(Z\in\mathcal C_m^G\) is zero if and only if
\[
 T(i)\wedge\Phi^H Z=0
 \qquad(H\le G,\ 0\le i\le m).
\]
For the image of a finite object, its support is tested equally by \(T(i)\) and \(K(i)\) at these heights.
For every compact object of \(\mathcal C_m^G\), membership in the support at
\((H,i)\) is equivalent to nonvanishing of \(K(i)\wedge\Phi^H(-)\).
\end{proposition}
\begin{proof}
Inflation of a finite type-\((m+1)\) spectrum generates the kernel as a localizing tensor ideal. Equivalently, use its products with the finite orbit generators. This is a finite localization, and its compact category is the idempotent completion of the corresponding compact quotient. The point classification and localization theorem of Balmer--Sanders identify its spectrum with the displayed subspace \cite[Theorems~4.9 and 4.14, Remark~5.3, and Proposition~5.24]{BS}. Indeed the inflated generator belongs precisely to the primes at test heights at most \(m\).

Geometric fixed points commute with this inflated localization. Each \(\Phi^H Z\) is therefore \(L_m^f\)-local; telescope detection at the retained heights, followed by genuine geometric detection, proves the zero criterion. For a finite representative, \eqref{eq:finite-tests} gives agreement of the two support tests.
For the assertion about all compact objects, the kernel of
\(K(i)\wedge\Phi^H(-)\) on compacts is a tensor-triangular prime by the
Morava \(K\)-theory field property. Its inverse image in the original compact
category is \(\mathcal P_i^G(H)\), since \(K(i)\wedge L_m^fS\simeq K(i)\)
for \(i\le m\). The injectivity of the spectrum map for finite localization
therefore identifies this kernel with the point labelled \((H,i)\).
\end{proof}
